% E7-E9, E12 — robustness battery for the hazard-guided controller. All ZERO-GPU,
% built on the frozen E6 precompute cache via experiments/e6_offline.py.
% Figures: results/figures/{e7_baselines,e8_ablation,e9_stats,e12_threshold}.png
% Tables : results/tables/{e7_baselines,e8_ablation,e9_stats,e12_threshold}.json
%
% HONESTY CONTRACT (do not soften back into overclaims):
%   * controller MATCHES always-on (not "beats"); hazard-always is within noise.
%   * detector edge over random is BORDERLINE and gsm8k-driven; bbh shows nothing.
%   * the significant result is rescue >> harm asymmetry.
%   * effective n = 184 questions (seed-collapse in the per-qid cache).

\subsection{Is the controller real? Baselines and ablation}
\label{sec:rq7}

A hazard-guided controller is only worth its complexity if it beats trivial policies
and if its gain is a genuine per-task signal rather than an artifact. We test both on
the same $n=516$ offline population, comparing every policy at a \emph{matched}
intervention budget (apply the safeguard to the top-$r$ debates by risk score, ties
broken at random).

\paragraph{Baselines (E7).}
Figure~\ref{fig:e7} pits the learned hazard against five deployable heuristics. The
round-0 \emph{disagreement} geometry is a strong baseline: intervening where agents
disagree, where answer-entropy is high, or where the majority margin is small all reach
$\approx\!61.3\%$ accuracy at a $40\%$ budget, versus $60.1\%$ for random. The learned
hazard is the \emph{best} single signal --- highest at every budget ($62.2\%$ at $40\%$)
and, more tellingly, the best ranker of true extinctions (AUPRC $0.375$ vs.\ $\le0.33$
for any heuristic and $0.225$ for random) --- but the margin over the structural
heuristics is small in the operating band and widens only at higher budget. Two
negative controls sharpen the story: an \emph{influence-only} detector (the F3
PageRank/stub\-bornness block) \emph{underperforms} the structural heuristics
(AUPRC $0.300$), and a \emph{revision} trigger is not even evaluable at the $g=0$ gate
--- no agent has revised yet --- so it coincides with random. The controller's power is
round-0 behavioral geometry, not influence and not revision dynamics.

\paragraph{Ablation (E8).}
Figure~\ref{fig:e8} decomposes the detector. Pooled (Fig.~\ref{fig:e8}A), no single
feature group dominates and \emph{influence is not required}: dropping the F3 block
actually raises AUPRC ($0.375\!\to\!0.415$), while text-surface features (F4) carry the
most unique pooled signal. But the pooled view is confounded --- GSM8K and BBH differ in
both extinction prevalence and verbosity --- so we validate \emph{within} each benchmark
(Fig.~\ref{fig:e8}B). Here the honest picture emerges: the detector \emph{ranks}
extinctions within GSM8K (the full detector attains AUPRC lift $1.98\times$; influence
alone $2.46\times$), yet \emph{targeting} those debates yields essentially no accuracy
edge over a random budget at a $40\%$ intervention rate --- $+0.0$ points within GSM8K,
$-1.2$ points within BBH, and only $+2.1$ points pooled. (An earlier per-question cache
that shared one outcome across a question's five seeds had inflated the GSM8K figure;
the numbers here use the corrected per-seed outcomes.) Ranking ability does not convert
into a controller advantage because the safeguard's rescue rate is modest and it harms
would-live debates (\S\ref{sec:rq9}). Influence remains benchmark-dependent (best ranker
in GSM8K, $\approx$random in BBH). For contrast, the pooled multi-observation-time E3a
detector, where the F2 dynamics block \emph{is} observable, improves monotonically
(AUPRC $0.19\!\to\!0.22\!\to\!0.24$ as dynamics then influence are added); dynamics help
\emph{when} available, but the controller's early gate cannot see them. The robust
reading: combine all channels for ranking, but the controller currently matches --- does
not beat --- a random budget on accuracy in either benchmark.

\subsection{Statistical robustness}
\label{sec:rq9}

Each debate now carries its own \emph{per-seed} counterfactual outcome (an earlier cache
that shared one outcome across a question's five seeds has been corrected). The five
seeds of a question are still correlated repeats of the same question, so the resampling
unit is the question, not the debate: the $516$ debates form $184$ clusters. We report a
cluster-bootstrap (resampling questions, $B=4000$, the deployed detector held fixed);
Figure~\ref{fig:e9} shows the intervals. Three conclusions follow. (i)~The controller
\emph{matches} always-on: the hazard$-$always accuracy gap at the peak operating point is
$+0.0$~points, CI $[-1.6,+1.2]$ --- squarely within noise. (ii)~The detector's edge over
a matched-random budget is \emph{suggestive but not conclusive}: the area under the
policy curve in the $20$--$40\%$ band is $+0.020$, CI $[-0.003,+0.042]$; every per-budget
interval includes zero, pooled and within each benchmark (no benchmark reaches
significance individually). (iii)~The one robustly significant effect is the
\emph{asymmetry} of intervention: applied everywhere, the safeguard rescues $39.5\%$ of
would-die debates (CI $[32.4,46.8]$) while harming $20.6\%$ of would-live debates
(CI $[15.5,26.1]$) --- non-overlapping intervals (and likewise within GSM8K). The
\emph{principle} of selective intervention is thus on firm ground; the current detector
ranks extinctions but does not yet convert that into a significant accuracy edge over a
random budget --- an honest gap for the mechanism, not the principle.

\subsection{Threshold and budget robustness}
\label{sec:rq12}

Finally we rule out the natural objection that $\tau^\star=0.52$ was chosen on the test
set (Figure~\ref{fig:e12}). Our recommended headline is not a threshold at all but the
\emph{budget-conditioned frontier} (Fig.~\ref{fig:e12}A): ``given an intervention budget
$r$, apply the safeguard to the $r$ highest-hazard debates.'' A budget is fixed a priori,
so nothing is tuned on test, and the frontier itself is the deliverable. To test whether
a \emph{data-chosen} threshold would generalize, we cross-fit it (GroupKFold-$5$ over
questions; on each training split pick $\tau^\star=\arg\max_\tau$ train accuracy and apply
it to the held-out split). The selected threshold is unstable across folds
($\tau^\star\in[0.004,0.52]$) and the pooled held-out accuracy is $61.4\%$,
CI $[56.2,66.7]$, at a mean $67\%$ intervention rate --- below the in-sample peak
($63.2\%$) and indistinguishable from always-on ($62.8\%$). The peak is therefore mildly
selection-optimistic, and a properly cross-validated policy simply matches always-on ---
exactly consistent with \S\ref{sec:rq9}.

\paragraph{Limitations and next steps.}
Three honest limitations bound these claims. \emph{First}, on aggregate accuracy the
controller \emph{matches} both always-on and a matched-random budget --- the edge over
random is not significant on either benchmark. Its demonstrated value is doing the job of
always-on at roughly half the interventions while avoiding a statistically significant
collateral harm; the accuracy payoff of \emph{targeting} (over random targeting) is a
real, open gap. \emph{Second}, only two benchmarks are covered here; a separate
detectability sweep (E11) shows extinction is even more predictable on StrategyQA
(hazard ROC $0.87$) than GSM8K, so a benchmark expansion is the natural next test of
whether targeting ever converts into accuracy. \emph{Third}, the three-agent ensemble is
fixed; whether the phenomenon survives other model families and $N>3$ is untested (E10).
(The per-question seed-collapse present in an earlier draft has been fixed here: every
seed now carries its own counterfactual outcome, which is what shrank the apparent
always-on headroom from $8.5$ to $4.5$ points and removed a spurious GSM8K advantage.) We
regard the mechanism --- predict-then-target, justified by the significant rescue/harm
asymmetry and a detector that ranks extinctions at $\sim\!2\times$ prevalence --- as the
contribution, and converting that ranking into a significant accuracy edge as the main
open problem, which \S\ref{sec:rq13} localizes precisely: the treatment \emph{benefit},
unlike the failure, is not predictable from round-0 features.

% ---------------------------------------------------------------------------
\begin{figure}[t]
  \centering
  \includegraphics[width=0.82\linewidth]{results/figures/e7_baselines.png}
  \caption{\textbf{E7 --- controller baselines at matched budget.} The learned hazard is
  the best single signal at every budget and the best extinction ranker (AUPRC $0.375$),
  but round-0 disagreement/entropy/margin heuristics are strong; \emph{influence-only}
  trails them and a \emph{revision} trigger is degenerate at the $g=0$ gate
  ($\equiv$ random). Bands are within-run random CIs.}
  \label{fig:e7}
\end{figure}

\begin{figure}[t]
  \centering
  \includegraphics[width=\linewidth]{results/figures/e8_ablation.png}
  \caption{\textbf{E8 --- ablation.} \textbf{(A)} Pooled, no group dominates and influence
  (F3) is not required (dropping it raises AUPRC). \textbf{(B)} Within-benchmark accuracy
  edge over a random budget at a $40\%$ intervention rate is near zero everywhere
  ($+0.0$ on GSM8K, $-1.2$ on BBH, $+2.1$ pooled): the detector ranks extinctions but
  targeting them barely beats random. Influence is benchmark-dependent (best on GSM8K,
  $\approx$random on BBH).}
  \label{fig:e8}
\end{figure}

\begin{figure}[t]
  \centering
  \includegraphics[width=\linewidth]{results/figures/e9_stats.png}
  \caption{\textbf{E9 --- cluster-bootstrap over $184$ questions.} \textbf{(A)}
  hazard$-$random is borderline (every interval includes zero, pooled and within each
  benchmark); hazard$-$always is within noise ($+0.0$ points). \textbf{(B)} the
  rescue$\gg$harm asymmetry is significant (non-overlapping intervals, pooled and GSM8K).
  The principle is solid; converting the detector's ranking into an accuracy edge is the
  open gap.}
  \label{fig:e9}
\end{figure}

\begin{figure}[t]
  \centering
  \includegraphics[width=\linewidth]{results/figures/e12_threshold.png}
  \caption{\textbf{E12 --- threshold/budget robustness.} \textbf{(A)} We report the
  budget-conditioned frontier (top-$r$ by hazard), which needs no test-time tuning.
  \textbf{(B)} A cross-fitted, data-chosen threshold yields $61.4\%$ held-out accuracy ---
  below the $63.2\%$ in-sample peak and $\approx$ always-on ($62.8\%$): the peak is mildly
  selection-optimistic, no cherry-pick advantage survives.}
  \label{fig:e12}
\end{figure}
